\documentclass[10pt,twocolumn]{article}
\usepackage[margin=0.75in,top=0.75in,bottom=0.75in]{geometry}
\usepackage{times}
\usepackage{graphicx}
\usepackage{booktabs}
\usepackage{amsmath}
\usepackage{hyperref}
\usepackage{enumitem}
\usepackage{xcolor}
\usepackage{caption}
\usepackage{subcaption}
\usepackage{microtype}
\usepackage{parskip}

\setlength{\parskip}{3pt}
\setlength{\parindent}{0pt}
\captionsetup{font=small, labelfont=bf}

\hypersetup{colorlinks=true, linkcolor=blue, citecolor=blue, urlcolor=blue}

\title{\textbf{High Alert: Behavioral Stage Discovery in Addiction\\
Treatment Reviews Using Clustering and LLM Reasoning}}
\author{Ben Blake\\
MSCS --- University of Missouri--Kansas City\\
\texttt{ben.blake@umkc.edu}\\
\small{\url{https://github.com/ben-blake/high-alert}}}
\date{}

\begin{document}

\maketitle
\thispagestyle{empty}

%-----------------------------------------------------------------
\section*{Abstract}
%-----------------------------------------------------------------
We present \textit{High Alert}, a public health analytics pipeline that
discovers recovery trajectory stages from patient reviews of addiction
treatment medications. Using UMAP dimensionality reduction, HDBSCAN
clustering, and LLM-based stage labeling grounded in the Transtheoretical
Model (TTM), we identify 10 behaviorally distinct patient populations from
7,192 filtered reviews spanning 2008--2017. A temporal spike detector
flags four quarters of elevated high-risk language, and LLM-generated
summaries translate each cluster into plain-language briefings for public
health analysts. All outputs are reproducible from a single shell command.

%-----------------------------------------------------------------
\section{Introduction}
%-----------------------------------------------------------------
Addiction treatment medication reviews contain rich, unsolicited signals
about patient recovery trajectories---side effects that threaten adherence,
sustained sobriety milestones, and relapse events. Mining these signals at
scale can reveal population-level patterns that clinical encounter data
misses: patients who relapse rarely return for follow-up appointments, but
they do write reviews.

Existing work on social media health mining focuses largely on opioid
mentions on Twitter or Reddit, using keyword search or supervised
classifiers that require labeled data. We take an unsupervised approach:
embed reviews, cluster them semantically, and have an LLM \textit{label}
the clusters---mapping each to a TTM stage and risk level without any
hand-labeled training examples.

Our contributions are: (1) an end-to-end reproducible pipeline from raw
CSV to analyst-ready summaries; (2) an empirical finding about survivorship
bias in drug review data; and (3) temporal spike detection that surfaces
quarters of elevated distress language aligned with known public health
events.

%-----------------------------------------------------------------
\section{Dataset}
%-----------------------------------------------------------------
We use the \textbf{UCI Drug Reviews} dataset (KUC Hackathon Winter
2018)~\cite{grasser2018}, which contains 215,063 patient reviews of
prescription medications sourced from \texttt{drugs.com}. Each row
includes drug name, medical condition, free-text review, numerical rating
(1--10), date, and a community usefulness count.

We filter to reviews matching 8 addiction-related conditions (e.g.,
\textit{Opiate Dependence}, \textit{Smoking Cessation}, \textit{Alcohol
Withdrawal}) and 15 addiction treatment drugs (e.g., Suboxone,
Varenicline, Naltrexone, Methadone, Nicotine replacement), yielding
\textbf{7,192 reviews} across 67 unique drugs and 27 conditions spanning
\textbf{2008-Q1 to 2017-Q4}. No personally identifiable information is
present; all data is public.

\begin{table}[h]
\centering
\caption{Corpus composition after filtering.}
\small
\begin{tabular}{lrr}
\toprule
\textbf{Top Conditions} & \textbf{N} & \textbf{\%} \\
\midrule
Smoking Cessation   & 2{,}440 & 33.9\% \\
Opiate Dependence   & 1{,}477 & 20.5\% \\
Alcohol Dependence  &   518 &  7.2\% \\
Opiate Withdrawal   &   434 &  6.0\% \\
Alcohol Withdrawal  &   285 &  4.0\% \\
\bottomrule
\end{tabular}
\label{tab:corpus}
\end{table}

%-----------------------------------------------------------------
\section{Data Preprocessing}
%-----------------------------------------------------------------
\textbf{Text cleaning.} Reviews are HTML-decoded, tag-stripped, and
lowercased using Python's \texttt{html} and \texttt{re} modules.

\textbf{Temporal features.} The date string (\texttt{27-Nov-16} format) is
parsed to extract year, quarter, and a \texttt{year\_quarter} key used for
all time-series analysis.

\textbf{Embeddings.} Cleaned reviews are encoded with
\textbf{all-MiniLM-L6-v2}~\cite{reimers2019} via
\texttt{sentence-transformers}, producing 384-dimensional cosine-normalized
vectors. Embeddings are cached to disk; the cache is reused when the corpus
is unchanged.

\textbf{Risk baseline.} A keyword classifier labels each review as
\textsc{high}, \textsc{moderate}, or \textsc{low} risk using 19 high-risk
patterns (e.g., \textit{relapsed}, \textit{overdose}, \textit{can't stop})
and 18 low-risk patterns (e.g., \textit{years sober}, \textit{smoke free}).
This produces: HIGH = 446 (6.2\%), MODERATE = 4,908 (68.2\%),
LOW = 1,838 (25.6\%).

%-----------------------------------------------------------------
\section{Methods}
%-----------------------------------------------------------------

\subsection{Dimensionality Reduction and Clustering}
Embeddings are reduced to 5 dimensions with
\textbf{UMAP}~\cite{mcinnes2018} (\texttt{n\_neighbors=15},
\texttt{min\_dist=0.1}, cosine metric), then clustered with
\textbf{HDBSCAN}~\cite{campello2013} (\texttt{min\_cluster\_size=130},
\texttt{min\_samples=10}). A separate 2-D UMAP projection is computed for
visualization. UMAP and KMeans use fixed random seeds for reproducibility.

\subsection{LLM Cluster Labeling}
For each cluster we select the $k=5$ reviews closest to the cluster
centroid (cosine similarity in 384-D space) and send them---along with
cluster statistics (dominant drugs, conditions, median rating) and a list
of all previously assigned stage names---to \textbf{llama-3.1-8b-instant}
via the Groq API. The model assigns a substance-qualified stage name
(e.g., \texttt{SUBOXONE\_LONG\_TERM\_SUCCESS}), a TTM stage, a one-sentence
description, and a risk level. A Python deduplication step guarantees all
stage names are unique, appending numeric suffixes when the LLM repeats a
name.

\subsection{LLM Risk Classification}
For evaluation, a stratified sample of 300 reviews is classified
individually by the LLM (batched 5 per call). Each review receives a stage
label, risk level, and one-sentence rationale. These LLM labels serve as
the evaluation reference against which we compare our baseline and
cluster-based risk assignments.

\subsection{Temporal Spike Detection}
For each calendar quarter we count baseline-flagged HIGH-risk reviews.
Global z-scores of the smoothed series (rolling window $w=4$) flag quarters
exceeding $z > 2.0$ as spikes. For each spike quarter, the LLM generates a
3-sentence public health narrative from the top-20 most useful HIGH-risk
reviews in that quarter.

\subsection{Explainability}
For each cluster, the LLM generates a structured 4-sentence summary: who
these patients are, their primary treatment challenges or successes, the
public health significance of the segment, and one recommended action for
health system planners.

%-----------------------------------------------------------------
\section{Experimental Design}
%-----------------------------------------------------------------
We compare three risk-classification approaches against LLM labels on 300
held-out reviews:

\begin{itemize}[noitemsep,topsep=2pt]
  \item \textbf{Baseline:} keyword/regex rules (37 patterns, no ML)
  \item \textbf{Cluster:} HDBSCAN cluster risk level assigned to each
        member review
  \item \textbf{LLM:} individual review classification (reference labels)
\end{itemize}

Macro-averaged F1 is reported across three risk classes. All hyperparameters
live in \texttt{config.yaml}; the full pipeline runs end-to-end with
\texttt{bash reproduce.sh}.

%-----------------------------------------------------------------
\section{Results and Discussion}
%-----------------------------------------------------------------

\subsection{Cluster Discovery}
HDBSCAN identified \textbf{10 behavioral clusters} with 389 noise points
(5.4\%). Figure~\ref{fig:umap} shows the 2-D UMAP projection. Clusters
separate cleanly by substance: Chantix/Varenicline smokers form a dense
cloud on the left; opioid-treatment patients (Suboxone, Vivitrol) occupy
distinct regions; the Librium alcohol-withdrawal cluster is small but
spatially isolated. All 10 clusters receive LOW risk labels from the LLM,
which is explained by the \textbf{survivorship bias} discussed below.

\begin{figure}[h]
  \centering
  \includegraphics[width=\linewidth]{umap_clusters.png}
  \caption{UMAP 2-D projection of 7,192 reviews colored by behavioral cluster.}
  \label{fig:umap}
\end{figure}

\subsection{Survivorship Bias Finding}
A key empirical finding is that HIGH-risk reviews (n=446) do not form
clusters---they appear as HDBSCAN noise. This is consistent with
\textit{survivorship bias}: patients in active relapse, overdose, or crisis
are unlikely to write medication reviews. Each relapse narrative is
idiosyncratic (different drug, circumstance, duration), while each success
narrative follows a similar arc, producing tight semantic clusters. The
high-risk population is best captured by the keyword baseline, not by
cluster membership. This finding has direct methodological implications:
embedding-based clustering should be used for characterizing stable
behavioral populations, while keyword signals remain essential for detecting
acute risk.

\subsection{Evaluation}
\begin{table}[h]
\centering
\caption{Macro-averaged F1 against LLM reference labels (n=300).}
\small
\begin{tabular}{lccc}
\toprule
\textbf{Approach} & \textbf{HIGH F1} & \textbf{LOW F1} & \textbf{Macro F1} \\
\midrule
Keyword Baseline & 0.021 & 0.036 & 0.288 \\
Cluster-Based    & 0.000 & 0.027 & 0.043 \\
\bottomrule
\end{tabular}
\label{tab:eval}
\end{table}

The keyword baseline outperforms the cluster approach on macro F1 (0.288
vs.\ 0.043), driven by high MODERATE-class recall. Neither approach
captures HIGH-risk reviews well because the LLM reference labels on 300
sampled reviews are also predominantly MODERATE (per-sample LLM labels
skew toward caution with limited context). The cluster approach assigns
risk at the cluster level, not the review level, and all clusters are LOW
risk---making HIGH-risk detection structurally impossible for this method.
These results confirm that keyword and embedding approaches are
\textit{complementary}: keywords find isolated high-risk signals;
clustering reveals the dominant behavioral populations.

\subsection{Temporal Analysis}
Figure~\ref{fig:spikes} shows quarterly HIGH-risk review counts.
Four spike quarters were detected: \textbf{2015-Q4}, \textbf{2016-Q1},
\textbf{2017-Q1}, and \textbf{2017-Q3}. The 2016-Q1 spike coincides with
growing clinical awareness of benzo-opioid co-prescription risks and rising
naloxone demand. The 2017 spikes align with U.S. opioid epidemic peak
reporting. LLM narratives for each spike describe common themes: tolerance
and rapid relapse, pain-addiction comorbidities, and the difficulty of
medication tapering.

\begin{figure}[h]
  \centering
  \includegraphics[width=\linewidth]{spike_detection.png}
  \caption{Quarterly HIGH-risk review count with detected spike quarters annotated.}
  \label{fig:spikes}
\end{figure}

%-----------------------------------------------------------------
\section{Ethical Considerations}
%-----------------------------------------------------------------
\textbf{No PII.} All data is public and self-reported to a public platform.
No names, dates of birth, or identifiers are present or inferred.

\textbf{Population-level only.} All outputs describe group-level patterns,
not individual risk scores. The system is designed for public health
planning, not clinical triage or surveillance of individuals.

\textbf{Framing.} Clusters and summaries are framed around treatment
support needs, not criminalization or personal failure. HIGH-risk language
detection is presented as a signal for resource allocation, not a patient
flag.

\textbf{Model limitations.} LLM labels are probabilistic and may reflect
biases in the training data. All LLM outputs are clearly marked as
AI-generated and should be reviewed by domain experts before informing
policy.

\textbf{Data access.} The raw dataset is publicly available on Kaggle. No
scraping, de-anonymization, or secondary collection was performed.

%-----------------------------------------------------------------
\section{Conclusion and Future Work}
%-----------------------------------------------------------------
We demonstrated that unsupervised semantic clustering combined with
LLM-based stage labeling can automatically discover interpretable recovery
trajectory segments from patient-generated text, without any labeled
training data. The pipeline is fully config-driven, reproducible, and
outputs analyst-ready markdown summaries and figures.

The survivorship bias finding is the most actionable result: high-risk
patients are underrepresented in review data, and this gap should inform
any public health study relying on voluntary self-report. Keyword-based
detection remains necessary to surface these hidden signals.

\textbf{Future work} includes: (1) integrating Reddit/social media data
where high-risk voices are more prevalent; (2) applying larger LLMs
(e.g., llama-3.1-70b) for richer cluster characterization; (3) training a
lightweight supervised classifier on LLM-labeled reviews to enable
real-time scoring; and (4) extending temporal analysis to detect emerging
drug trends before they appear in clinical data.

%-----------------------------------------------------------------
\begin{thebibliography}{9}
\bibitem{grasser2018}
F.\ Gr\"{a}{\ss}er et al., ``Aspect-Based Sentiment Analysis of Drug Reviews,''
\textit{Proc.\ 2018 Int.\ Conf.\ on Digital Health}, 2018.

\bibitem{reimers2019}
N.\ Reimers and I.\ Gurevych, ``Sentence-BERT: Sentence Embeddings using
Siamese BERT-Networks,'' \textit{EMNLP}, 2019.

\bibitem{mcinnes2018}
L.\ McInnes, J.\ Healy, and J.\ Melville, ``UMAP: Uniform Manifold Approximation
and Projection for Dimension Reduction,'' \textit{arXiv:1802.03426}, 2018.

\bibitem{campello2013}
R.\ Campello, D.\ Moulavi, and J.\ Sander, ``Density-Based Clustering Based on
Hierarchical Density Estimates,'' \textit{PAKDD}, 2013.
\end{thebibliography}

\end{document}
